% ===========================================================================
% fig-frame.tex -- v2. The residual frame, the encoding fork, and the SCORED
% object. \input by Sections/04c-reencoding.tex. Compiles against
% thesis_v2/preamble.tex and nothing else; no external data, no
% \includegraphics. FULLWIDTH, and MEASURED rather than assumed: the
% tikzpicture's own bounding box is 491.31pt wide against a fullwidth measure
% of 495.08pt (\textwidth 404.03pt + \marginparwidth + \marginparsep). The
% measurement is made by the picture itself, guarded by \fframedebug;
% figs/_harness_frame.tex switches it on, and it is inert in main.tex.
%
% RULE INHERITED FROM preamble.tex section 8: a page carrying a fullwidth float
% carries NO \marginpar. Every note in this figure is therefore drawn INSIDE
% the tikzpicture. Nothing here calls \gloss, \seealso or any margin wrapper.
%
% EVERY NUMBER IN THIS FIGURE, WITH ITS SOURCE. Nothing is invented.
%   34.6 / 200  (17%)  full profile, cell-line scope
%   118.2 / 200 (59%)  residual,     cell-line scope
%   120.3              residual,     plate scope            -- caption only
%   0.742 -> 0.743     within-well split-half retrieval reference, full ->
%                      residual (stored 0.741986901358855 and 0.7432424792)
%     RESULTS_cluster/target_divergence.json, scopes.cellline.{full,residual}
%     .divergence_mean and .ceiling_retrieval, and scopes.plate.residual;
%     and tab:divergence, Sections/04c-reencoding.tex lines 42-44.
%   100 genes per block, [DOWN], [END_CELL], weight 1/log2(rank+1), cosine,
%   NIR over other drugs in the same cell line
%     Sections/04c-reencoding.tex, "How the residual becomes a string, and how
%     it is scored"; and endcell/analysis/residual_eval.py,
%     signed_rank_from_sentence / signed_rank_from_vector (k = 100,
%     v[gi] = sgn / np.log2(r + 1.0)).
%   946 panel genes -- Sections/04c-reencoding.tex, "a 946-dimensional profile".
%   the stem weights: 1/log2(r+1) for r = 1..8 =
%     1.0000, 0.6309, 0.5000, 0.4307, 0.3869, 0.3562, 0.3333, 0.3155.
%     Every stem in stage 3 is one of these eight numbers; the +1 / -1 ticks
%     are there so a reader can check them.
%   "fewer than three other training drugs yields no generic and the condition
%     is dropped" -- Sections/04c-reencoding.tex, "Second, the scope is the
%     plate".
%   [DOWN] registered atomic / unregistered splits into sub-words --
%     Sections/03-Apparatus.tex line 283, and 00-HowToRead.tex line 209.
%
% WHAT IS ILLUSTRATIVE, AND SAYS SO IN THE CAPTION:
%   * the six gene profiles in stage 1. No axis is drawn; they carry no
%     measurement. Their ARITHMETIC is exact bar by bar, which is the only
%     thing a reader may take from them:
%       treated  {17, 9,14,12, 8, 6,11, 5}
%       control  {10,10, 8, 8, 7, 7,10, 6}
%       shift    { 7,-1, 6, 4, 1,-1, 1,-1}  = treated - control
%       generic  { 3, 2, 3, 2,-2, 2,-2, 2}
%       residual { 4,-3, 3, 2, 3,-3, 3,-3}  = shift - generic
%     Chosen so the residual carries the MAJORITY of the shift's squared norm
%     (74 of 106 here), because fig:residual-energy measures 62.1%
%     drug-specific and a figure drawing the residual as a sliver would
%     contradict it. Both figures are UNCENTRED sums of squares, each a
%     squared L2 norm and NOT a variance: the shift {7,-1,6,4,1,-1,1,-1} has
%     mean 2, not 0, and 49+1+36+16+1+1+1+1 = 106; the residual gives 74 the
%     same way. The label fig:residual-energy and the file
%     figs/04c-residual-energy.tex keep the older spelling because they are a
%     cross-reference target and an \input target.
%     v1 drew {1,0,0,1,0,1,0,0}: five of its eight bars had height zero and
%     rendered as nothing at all, so the one object the thesis is about read
%     as "missing" or "to be determined". FIXED HERE, and drawn SIGNED,
%     because the sign is what splits the two blocks in stage 2.
%   * which gene sits at which coordinate in stage 3, and which two the
%     prediction puts on the wrong side of the boundary. The STEM HEIGHTS are
%     not illustrative -- they are the exact weights.
%   * no cosine VALUE is printed anywhere. The two stem vectors are drawn, not
%     measured, so a number there would be an invented result.
%
% LAYOUT NOTES, all forced by a build:
%   * three stages, two dividers, one header rule, one brace. The brace spans
%     all THREE because the retrieval reference is an end-of-pipeline quantity:
%     it is a retrieval score computed in the scored space of stage 3, so
%     stopping it at stage 2 would misattribute it.
%   * stage 2's target string is ONE line. Wrapped over two it read as two
%     separate sequences and lost the up -> [DOWN] -> down ordering, which is
%     the only reason [DOWN] is drawn at all. Every chip is a NAMED node
%     chained off its predecessor's east anchor, so the two braces attach to
%     real node boundaries rather than to hand-computed x's. That was not
%     fastidiousness: the hand estimate of the string's width was 6mm short of
%     the truth, and on the first build [END_CELL] crossed the divider into
%     stage 3.
%   * type sizes follow figs/circularity.tex: \footnotesize (10pt, the caption
%     size) for the sentences a reader must read, \scriptsize (8pt) for symbol
%     labels, chips, brace labels and subordinate notes. Nothing is \tiny.
%   * stage 1 labels are centred on profiles 1.85cm apart. At 1.45cm
%     "generic(c)" and "residual r(d,c)" touch (v1's note, still true). They
%     are \scriptsize and not \footnotesize for the same reason: at 10pt they
%     collide even at 1.85cm, which a build showed.
% ===========================================================================
\begin{figure}[htbp]
\begin{fullwidth}
\centering
\begin{tikzpicture}[
  x=1cm, y=1cm, font=\small,
  bar/.style={fill=ink!72, draw=none},
  faintbar/.style={fill=ink!30, draw=none},
  base/.style={rulegrey, line width=0.25pt},
  gene/.style={draw=rulegrey, rounded corners=1pt, inner sep=1.6pt,
               font=\ttfamily\scriptsize, text=ink, anchor=west},
  shared/.style={gene, fill=rulegrey!12, text=quietink, draw=rulegrey!55},
  differs/.style={gene, fill=rulegrey!32, draw=ink!75, thick},
  atomic/.style={rounded corners=1pt, inner sep=1.6pt, fill=ink, draw=ink,
                 font=\ttfamily\scriptsize, text=white, anchor=west},
  dots/.style={anchor=west, font=\scriptsize, text=quietink},
  lbl/.style={font=\scriptsize, text=quietink},
  strong/.style={font=\scriptsize, text=ink},
  said/.style={anchor=west, font=\footnotesize, text=ink},
  op/.style={font=\normalsize, text=quietink},
  stagehead/.style={anchor=west, font=\sffamily\footnotesize\bfseries,
                    text=accent},
  note/.style={anchor=north west, font=\scriptsize, text=quietink,
               align=flush left, inner sep=0pt},
]

% ---- one bar profile. #1 x, #2 baseline y, #3 signed heights, #4 style.
% Heights are in "units"; one unit is 0.052cm. Bars are 0.125cm wide on a
% 0.160cm pitch, so eight span 1.245cm. The baseline rule is drawn under EVERY
% profile, because three of the six are signed and a bar below the line has to
% read as negative rather than as a misprint.
\newcommand{\fframeprof}[4]{%
  \foreach \hh [count=\ii from 0] in {#3} {%
    \pgfmathsetmacro{\bxx}{#1 + \ii*0.160}%
    \pgfmathsetmacro{\byy}{\hh*0.052}%
    \fill[#4] (\bxx,#2) rectangle ++(0.125,\byy);%
  }%
  \draw[base] (#1,#2) -- ++(1.245,0);%
}

% ---- one signed-rank stem plot. #1 zero-line y, #2 slot/weight list.
% Slots sit on a 0.255cm pitch from x = 13.16; a weight of 1 is 0.34cm.
\newcommand{\fframestems}[2]{%
  \draw[base] (13.10,#1) -- (17.02,#1);%
  \foreach \ss/\ww in {#2} {%
    \pgfmathsetmacro{\sxx}{13.16 + (\ss-1)*0.255}%
    \pgfmathsetmacro{\syy}{#1 + \ww*0.34}%
    \draw[ink!72, line width=0.55pt] (\sxx,#1) -- (\sxx,\syy);%
    \fill[ink!72] (\sxx,\syy) circle (0.5pt);%
  }%
  \draw[rulegrey, line width=0.3pt] (12.96,{#1-0.34}) -- (12.96,{#1+0.34});%
  \draw[rulegrey, line width=0.3pt] (12.92,{#1+0.34}) -- (13.00,{#1+0.34});%
  \draw[rulegrey, line width=0.3pt] (12.92,{#1-0.34}) -- (13.00,{#1-0.34});%
  \node[anchor=east, font=\scriptsize, text=quietink]
       at (12.88,{#1+0.34}) {$+1$};%
  \node[anchor=east, font=\scriptsize, text=quietink]
       at (12.88,{#1-0.34}) {$-1$};%
}

% ================================================== the header band
\node[stagehead] at (0.05,3.80) {1\,\textperiodcentered\,the residual is built};
\node[stagehead] at (5.70,3.80) {2\,\textperiodcentered\,it becomes a string};
\node[stagehead] at (12.60,3.80) {3\,\textperiodcentered\,it becomes a vector};
\draw[rulegrey, line width=0.4pt] (0.05,3.62) -- (17.30,3.62);
\draw[rulegrey!70, line width=0.3pt] (5.50,-1.66) -- (5.50,3.54);
\draw[rulegrey!70, line width=0.3pt] (12.40,-1.66) -- (12.40,3.54);

% ================================================== STAGE 1: two subtractions
\node[lbl]    at (0.67,3.36) {treated};
\node[lbl]    at (2.52,3.36) {control};
\node[strong] at (4.37,3.36) {shift $\shift(d,c)$};
\fframeprof{0.05}{2.30}{17,9,14,12,8,6,11,5}{bar}
\node[op] at (1.53,2.58) {$-$};
\fframeprof{1.90}{2.30}{10,10,8,8,7,7,10,6}{bar}
\node[op] at (3.38,2.58) {$=$};
\fframeprof{3.75}{2.30}{7,-1,6,4,1,-1,1,-1}{bar}

\node[strong] at (0.67,1.90) {shift};
\node[lbl]    at (2.52,1.90) {generic$(c)$};
\node[strong] at (4.37,1.90) {residual $\resid(d,c)$};
\fframeprof{0.05}{1.30}{7,-1,6,4,1,-1,1,-1}{bar}
\node[op] at (1.53,1.42) {$-$};
\fframeprof{1.90}{1.30}{3,2,3,2,-2,2,-2,2}{faintbar}
\node[op] at (3.38,1.42) {$=$};
\fframeprof{3.75}{1.30}{4,-3,3,2,3,-3,3,-3}{bar}

\node[note, text width=5.30cm] (nG) at (0.05,0.88)
  {generic$(c)$ is the mean over the \emph{other} training drugs in the same
   plate group; a group with fewer than three of them yields no generic and the
   condition is dropped. The residual is drawn signed because its sign is what
   splits the two blocks in stage~2.};

% ================================================== STAGE 2: the encoding fork
% Every gene label is a real panel gene, so every string drawn could occur in
% this data. The shared positions carry four of the most abundant panel genes by
% mean measured expression over the 606 consensus truth profiles
% (nir_consensus_profiles.npz: GAPDH rank 1, RPS6 2, HSPD1 7, ALDOA 8; GAPDH and
% RPS6 agree with gene_identity_check.json). They were ACTB, TPT1 and EEF1A1,
% none of which is in the 946-gene panel. In row (b) DDB2, an in-panel p53
% target, replaces MDM2, which is not in the panel. CDKN1A and HMOX1 are in the
% panel and unchanged. Labels only: no position, count or number changed.
\node[said] at (5.62,3.30) {\textbf{(a)} genes ranked by \emph{expression}};
\node[shared]  (a1) at (5.62,2.88) {GAPDH};
\node[shared]  (a2) at ([xshift=3pt]a1.east) {RPS6};
\node[shared]  (a3) at ([xshift=3pt]a2.east) {HSPD1};
\node[differs] (a4) at ([xshift=3pt]a3.east) {CDKN1A};
\node[shared]  (a5) at ([xshift=3pt]a4.east) {ALDOA};
\node[dots]    (a6) at ([xshift=4pt]a5.east) {$\cdots$};
\node[shared]  (b1) at (5.62,2.48) {GAPDH};
\node[shared]  (b2) at ([xshift=3pt]b1.east) {RPS6};
\node[shared]  (b3) at ([xshift=3pt]b2.east) {HSPD1};
\node[differs] (b4) at ([xshift=3pt]b3.east) {HMOX1};
\node[shared]  (b5) at ([xshift=3pt]b4.east) {ALDOA};
\node[dots]    (b6) at ([xshift=4pt]b5.east) {$\cdots$};
\node[said] at (5.62,2.06)
  {only \textbf{34.6 / 200} tokens differ \;(\textbf{17\%})};

\node[said] at (5.62,1.56)
  {\textbf{(b)} genes ranked by \emph{signed residual}};
\node[differs] (c1) at (5.62,1.14) {CDKN1A};
\node[differs] (c2) at ([xshift=3pt]c1.east) {DDB2};
\node[dots]    (c3) at ([xshift=3pt]c2.east) {$\cdots$};
\node[atomic]  (c4) at ([xshift=4pt]c3.east) {[DOWN]};
\node[differs] (c5) at ([xshift=4pt]c4.east) {HMOX1};
\node[dots]    (c6) at ([xshift=3pt]c5.east) {$\cdots$};
\node[atomic]  (c7) at ([xshift=4pt]c6.east) {[END\_CELL]};
\draw[decorate, decoration={brace, mirror, amplitude=3pt}, rulegrey]
  (c1.south west) ++(0,-0.06) -- ($(c3.south east)+(0,-0.06)$);
\node[lbl, anchor=north west] at ($(c1.south west)+(0,-0.16)$)
  {$100$ genes, $\resid > 0$};
\draw[decorate, decoration={brace, mirror, amplitude=3pt}, rulegrey]
  (c5.south west) ++(0,-0.06) -- ($(c6.south east)+(0,-0.06)$);
\node[lbl, anchor=north west] at ($(c5.south west)+(0,-0.16)$)
  {$100$ genes, $\resid < 0$};
\node[said] at (5.62,0.20)
  {\textbf{118.2 / 200} tokens differ \;(\textbf{59\%})};

\node[atomic] (dn) at (5.62,-0.28) {[DOWN]};
\node[note, text width=5.35cm] (nD) at ($(dn.east)+(0.14,0.11)$)
  {is \emph{one} registered atomic token. Unregistered it splits into
   sub-words, and the up\slash down boundary stops being a clean symbol.};

% ================================================== STAGE 3: the scored vector
\node[note, text width=4.70cm] (nW) at (12.60,3.44)
  {signed rank weights $w_r = 1/\log_2(r{+}1)$};

% the two coordinates on which the prediction and the truth carry OPPOSITE
% SIGNS -- the boundary error the [DOWN] token exists to prevent. Drawn FIRST
% so the stems sit on top of it.
\fill[rulegrey!18] (14.82,0.70) rectangle (15.33,2.42);
% The band's label sits in the GAP BETWEEN the two plots, running right from the
% band's right edge. Above the plots it ran into the second line of the weights
% note, which is math and taller than a text line; centred on the band it set
% flush against "true v" and the two read as one phrase. In the gap the nearest
% ink is the predicted plot's slot-10 stem, which stops at y = 1.72.
% anchor=EAST at 17.30, not west at 15.42: anchored west, this one 2.07cm label
% was the widest thing in the picture and pushed the bounding box out to
% 497.60pt, 1.35pt past the fullwidth measure -- one Overfull \hbox, which
% tools/build.sh caps at zero.
\node[lbl, anchor=east] at (17.30,1.46) {opposite signs};

\node[strong, anchor=west] at (13.16,2.50) {predicted $\hat{v}$};
\fframestems{2.06}{1/0.6309, 2/1.0000, 3/0.5000, 4/0.3869, 5/0.4307,
                   6/0.3562, 7/0.3155, 8/-0.6309, 9/0.3333, 10/-1.0000,
                   11/-0.5000, 12/-0.4307, 13/-0.3562, 14/-0.3869,
                   15/-0.3333, 16/-0.3155}

\node[strong, anchor=west] at (13.16,1.50) {true $v$};
\fframestems{1.06}{1/1.0000, 2/0.6309, 3/0.5000, 4/0.4307, 5/0.3869,
                   6/0.3562, 7/0.3333, 8/0.3155, 9/-1.0000, 10/-0.6309,
                   11/-0.5000, 12/-0.4307, 13/-0.3869, 14/-0.3562,
                   15/-0.3333, 16/-0.3155}

\node[note, text width=4.70cm] (nA) at (12.60,0.50)
  {$946$ coordinates, $200$ of them non-zero. $\cos(\hat{v},v)$, ranked
   against the other drugs in the same cell line, is $\NIR$.};

\node[note, text width=4.70cm] (nB) at (12.60,-0.60)
  {\DOWN{} has a token position in stage~2 but \emph{no coordinate} here: the
   sign carries the boundary.};

% ================================================== the invariant
\draw[decorate, decoration={brace, amplitude=4pt}, quietink, line width=0.5pt]
  (0.05,-1.80) -- (17.30,-1.80);
\node[anchor=north, font=\footnotesize, text=ink] (nBR) at (8.67,-1.98)
  {and the within-well split-half retrieval reference is \textbf{unchanged}:
   $0.742 \rightarrow 0.743$};

% ---- MEASUREMENT, left in and inert. \fframedebug is undefined in main.tex, so
% this expands to nothing there; figs/_harness_frame.tex \def's it and the
% picture then reports its own extent, the right edge of the target string, and
% the east edge of every node that could plausibly be the widest thing in the
% drawing. All three mattered: the string crossed a divider on the first build,
% and the bounding box ran 1.35pt past the fullwidth measure on a later one.
\ifx\fframedebug\undefined\else
  \path let \p1=(current bounding box.south west),
            \p2=(current bounding box.north east),
            \p3=(c7.east), \p4=(a6.east) in
    \pgfextra{\typeout{FRAME bbox x \x1 .. \x2 ; y \y1 .. \y2}%
              \typeout{FRAME target string ends \x3 ; row (a) ends \x4}};
  \path let \p5=(nW.east), \p6=(nA.east), \p7=(nB.east), \p8=(nBR.east),
            \p9=(nD.east), \p{10}=(nG.east) in
    \pgfextra{\typeout{FRAME east nW \x5 nA \x6 nB \x7 nBR \x8 nD \x9 nG \x{10}}};
\fi
\end{tikzpicture}
\caption[The residual, the target string and the scored vector]{\captiontitle{The
information was always present; the standard target's tokens barely expose it
--- and the string the model writes is not the vector that scores it.}
\emph{Stage 1.} The drug-specific residual is built by two subtractions. The
control removes the cell's own state, leaving the \emph{shift}; the mean over
the other drugs in the same plate group removes the \emph{generic} programme
that every drug triggers, leaving what makes this drug different. The six
gene profiles are illustrative shapes and carry no measurement --- no axis is
drawn --- but the arithmetic within the figure is exact bar by bar, and the
residual is drawn signed because its sign is what splits the two blocks in
stage~2. \emph{Stage 2.} That same measurement, written two ways. Ranking genes
by expression (a) leads every drug's profile with largely the same abundant
genes: only $34.6$ of the top $200$ tokens differ between two drugs in the same
context; ranking by signed residual
(b) raises that to $118.2$. (Divergence counts quoted at cell-line scope; Table~\ref{tab:divergence} quotes the plate scope this build trains in, $120.3$, and gives the cell-line figure in its note.) Row (b) is the
target string itself: a $100$-gene up block, \DOWN, a $100$-gene down block,
\ENDCELL. \emph{Stage 3.} What $\NIR$ ranks is not that string but a signed
vector over the $946$ panel genes, top-$k$ up positive and top-$k$ down
negative. The stem heights are the exact weights $1/\log_2(\textrm{rank}+1)$, so
the eight drawn per sign are $1.000$, $0.631$, $0.500$, $0.431$, $0.387$,
$0.356$, $0.333$ and $0.316$; which gene sits at which coordinate is
illustrative, and no cosine value is printed, because these two vectors are
drawn rather than measured. Stage~2 and stage~3 are different objects and only
stage~3 is scored. The bracket is the point of the figure --- the within-well
split-half retrieval reference is essentially identical under both encodings, so
no information has been added. What changed is how much of it the token sequence
exposes to the loss --- a property of the encoding, which is not the same as
showing the encoding is what binds (Table~\ref{tab:divergence}).}
\label{fig:frame}
\end{fullwidth}
\end{figure}
